\chapter{Tự khen trước một điều}
\begin{tinhhuongbox}{Tình huống | Situation}
\songngu{Nếu giờ chủ nhiệm mở bằng bảng lỗi, cả lớp sẽ co người lại rất nhanh.}{If homeroom opens with a board full of mistakes, the whole class will shrink quickly.}
\songngu{Khen trước một chút không làm buổi sinh hoạt yếu đi; nó làm lớp chịu hợp tác hơn.}{Starting with one piece of praise does not weaken the session; it makes the class more willing to cooperate.}
\end{tinhhuongbox}
\begin{noitrenlop}
\songngu{Trước khi sửa nhau, lớp mình tự khen một điều đi, cho công bằng với cuộc đời.}{Before we correct ourselves, let us praise ourselves for one thing, just to be fair to life.}
\end{noitrenlop}
\begin{vihieuqua}
\songngu{Câu này làm không khí mềm xuống mà không hề làm mất trọng tâm.}{This line softens the atmosphere without losing focus at all.}
\end{vihieuqua}
\begin{englishpocket}
\songngu{Cụm nên nhớ: \textit{Just to be fair to life.}}{Phrase to keep: \textit{Just to be fair to life.}}
\end{englishpocket}
\chapter{Bớt điều gì để lớp vui hơn}
\begin{tinhhuongbox}{Tình huống | Situation}
\songngu{Một lớp không chỉ cần thêm điều tốt, mà còn cần bớt điều làm không khí nặng đi.}{A class does not only need more good things; it also needs less of what makes the room heavier.}
\songngu{Câu hỏi bớt cái gì thường thực tế hơn câu hỏi đổi cả đời.}{Asking what to reduce is often more practical than asking how to change everything.}
\end{tinhhuongbox}
\begin{noitrenlop}
\songngu{Lớp mình cần bớt điều gì để tuần tới đi học đỡ giống leo dốc hơn?}{What does our class need less of so that coming to school next week feels less like climbing a hill?}
\end{noitrenlop}
\begin{vihieuqua}
\songngu{Lớp dễ nói thật với kiểu câu hỏi gần đời như thế này.}{Students are more likely to answer honestly when the question feels this close to life.}
\end{vihieuqua}
\begin{englishpocket}
\songngu{Cụm nên nhớ: \textit{Feels less like climbing a hill.}}{Phrase to keep: \textit{Feels less like climbing a hill.}}
\end{englishpocket}
\chapter{Một lời hứa nhỏ nhưng thật}
\begin{tinhhuongbox}{Tình huống | Situation}
\songngu{Lời hứa quá đẹp thường sống không lâu.}{Promises that sound too beautiful usually do not live very long.}
\songngu{Một điều nhỏ nhưng thật sẽ đáng tin hơn nhiều.}{One small but real promise is much more believable.}
\end{tinhhuongbox}
\begin{noitrenlop}
\songngu{Lớp hứa với nhau một điều ít thôi, nhưng làm thật, đừng hứa như đang phát biểu nhận giải.}{Let us promise each other just one small thing, but make it real. Do not promise like you are giving an award speech.}
\end{noitrenlop}
\begin{vihieuqua}
\songngu{Câu này kéo lớp về thực tế và làm cam kết bớt màu mè.}{This line pulls the class back to reality and makes the commitment less theatrical.}
\end{vihieuqua}
\begin{englishpocket}
\songngu{Cụm nên nhớ: \textit{Make it real.}}{Phrase to keep: \textit{Make it real.}}
\end{englishpocket}
\chapter{Một điều muốn cảm ơn}
\begin{tinhhuongbox}{Tình huống | Situation}
\songngu{Giờ chủ nhiệm chỉ sửa sai thì lớp sẽ mệt dần theo tuần.}{If homeroom only corrects mistakes, the class will grow tired week after week.}
\songngu{Một câu cảm ơn ngắn có thể cứu rất nhiều cảm giác tập thể.}{One short thank-you line can save a lot of group feeling.}
\end{tinhhuongbox}
\begin{noitrenlop}
\songngu{Trước khi kết tội nhau tiếp, ai có một điều muốn cảm ơn lớp, cảm ơn bạn, hay cảm ơn tổ mình không?}{Before we continue putting each other on trial, does anyone want to thank the class, a friend, or their group for something?}
\end{noitrenlop}
\begin{vihieuqua}
\songngu{Câu này vừa hài vừa làm không khí ấm lên rõ rệt.}{This line is funny and also makes the room noticeably warmer.}
\end{vihieuqua}
\begin{englishpocket}
\songngu{Cụm nên nhớ: \textit{Does anyone want to thank...?}}{Phrase to keep: \textit{Does anyone want to thank...?}}
\end{englishpocket}
\chapter{Tuần tới muốn được nhớ vì điều gì}
\begin{tinhhuongbox}{Tình huống | Situation}
\songngu{Một lớp mạnh là lớp biết tự tạo hình ảnh cho mình, không chỉ chờ giáo viên gắn nhãn.}{A strong class knows how to shape its own identity instead of waiting for the teacher to label it.}
\songngu{Câu hỏi này kéo giờ chủ nhiệm từ sửa lỗi sang xây lớp.}{This question moves homeroom from fixing mistakes to building a class.}
\end{tinhhuongbox}
\begin{noitrenlop}
\songngu{Tuần tới lớp mình muốn được nhớ vì điều gì, đừng để người ta nhớ vì ồn là mệt lắm.}{What do we want to be remembered for next week? Please do not let it be noise, because that is exhausting.}
\end{noitrenlop}
\begin{vihieuqua}
\songngu{Lớp bắt đầu nghĩ theo hướng danh dự tập thể chứ không chỉ theo hướng né phê bình.}{The class starts thinking in terms of shared pride, not only in terms of avoiding criticism.}
\end{vihieuqua}
\begin{englishpocket}
\songngu{Cụm nên nhớ: \textit{What do we want to be remembered for?}}{Phrase to keep: \textit{What do we want to be remembered for?}}
\end{englishpocket}